% GENERATED FILE. Edit canonical lessons, then run npm run generate:books.
% canonical-source-hash: fnv1a64:cfbd3f24cc8b034f
% canonical-lessons: TA-C514-tolircankam, TA-C514-velainiruttam, TA-C514-panimanai, TA-C514-kitanku, TA-C514-perunar

\chapter{Unions and Warehouses}
\label{ch:ta-unions-and-warehouses}
\hlchaptermodality{514}

\begin{chapteropening}
\textbf{By the end of this chapter:} \emph{I can talk about strikes, unions and storing goods.}

Complete the last lesson of chapter 514: I can talk about strikes, unions and storing goods.
\end{chapteropening}

\section[\texorpdfstring{\ta{தொழிற்சங்கம்} (toḻiṟcaṅkam)}{தொழிற்சங்கம் (toḻiṟcaṅkam)}]{\ta{தொழிற்சங்கம்} (toḻiṟcaṅkam) — a trade union}

\label{lesson:TA-C514-tolircankam}



\begin{quote}
Before the new one: say the Tamil for a weaver, then the Tamil for a maid.
\end{quote}

\subsection*{You'll want to know: \ta{தொழிற்சங்கம்}}
\textbf{\ta{தொழிற்சங்கம்}} — \emph{toḻiṟcaṅkam} — \textquotedblleft{}a trade union\textquotedblright{}.

\begin{quote}
\ta{தொழிற்சங்கம்} \ta{சம்பள} \ta{உயர்வு} \ta{கேட்டது}.
\end{quote}

\begin{quote}
\emph{toḻiṟcaṅkam campaḷa uyarvu kēṭṭatu.}
\end{quote}

\begin{quote}
The trade union asked for a pay rise.
\end{quote}

Before \ta{ச} the final -\ta{ல்} of \ta{தொழில்} becomes -\ta{ற்}: \ta{தொழில்} + \ta{சங்கம்} = \ta{தொழிற்சங்கம்}.

\begin{cousinweb}
\textbf{\ta{தொழிற்சங்கம்}} (\emph{toḻiṟcaṅkam}) joins the native \ta{தொழில்} (\emph{toḻil}), \textquotedblleft{}a trade\textquotedblright{} (DEDR 3524), and \ta{சங்கம்} (\emph{caṅkam}), \textquotedblleft{}an association\textquotedblright{}, from Sanskrit \emph{saṅgha} (Monier-Williams).
\end{cousinweb}

\begin{grammarlens}[title={What you've built}]
One more word for describing what you lived and felt.
\end{grammarlens}

\subsection*{Your turn}
\begin{itemize}
\raggedright
  \item \emph{Say it:} \emph{toḻiṟcaṅkam}
  \item \emph{Say it:} \emph{toḻiṟcaṅkam}, once more
  \item \emph{Say it:} \emph{toḻiṟcaṅkam campaḷa uyarvu kēṭṭatu.} — The trade union asked for a pay rise.
  \item \emph{Recall:} say the Tamil for a weaver, then the Tamil for a maid, then say \emph{toḻiṟcaṅkam} again
\end{itemize}

\begin{tcolorbox}[breakable,colback=teal!4,colframe=teal!35!black,title={Before you move on}]
What does \ta{தொழிற்சங்கம்} mean? (\textquotedblleft{}A trade union\textquotedblright{}.) Say it once more.
\end{tcolorbox}
\section[\texorpdfstring{\ta{வேலைநிறுத்தம்} (vēlainiṟuttam)}{வேலைநிறுத்தம் (vēlainiṟuttam)}]{\ta{வேலைநிறுத்தம்} (vēlainiṟuttam) — a strike}

\label{lesson:TA-C514-velainiruttam}



\begin{quote}
Before the new one: say the Tamil for a maid, then the Tamil for a trade union.
\end{quote}

\subsection*{You'll want to know: \ta{வேலைநிறுத்தம்}}
\textbf{\ta{வேலைநிறுத்தம்}} — \emph{vēlainiṟuttam} — \textquotedblleft{}a strike\textquotedblright{}.

\begin{quote}
\ta{வேலைநிறுத்தத்தால்} \ta{இன்று} \ta{பேருந்துகள்} \ta{ஓடவில்லை}.
\end{quote}

\begin{quote}
\emph{vēlainiṟuttattāl iṉṟu pēruntukaḷ ōṭavillai.}
\end{quote}

\begin{quote}
Because of the strike the buses did not run today.
\end{quote}

Word for word \textquotedblleft{}a stopping of work\textquotedblright{}: \ta{வேலை} + \ta{நிறுத்தம்}, from \ta{நிறுத்து} (\emph{niṟuttu}), to stop.

\begin{cousinweb}
\textbf{\ta{வேலைநிறுத்தம்}} (\emph{vēlainiṟuttam}) is a native compound: \ta{வேலை} (\emph{vēlai}), \textquotedblleft{}work\textquotedblright{} (DEDR 5540), and \ta{நிறுத்தம்} (\emph{niṟuttam}), \textquotedblleft{}a stopping\textquotedblright{}, from \ta{நில்} (\emph{nil}), \textquotedblleft{}to stand, to stop\textquotedblright{} (DEDR 3675).
\end{cousinweb}

\begin{grammarlens}[title={What you've built}]
One more word for describing what you lived and felt.
\end{grammarlens}

\subsection*{Your turn}
\begin{itemize}
\raggedright
  \item \emph{Say it:} \emph{vēlainiṟuttam}
  \item \emph{Say it:} \emph{vēlainiṟuttam}, once more
  \item \emph{Say it:} \emph{vēlainiṟuttattāl iṉṟu pēruntukaḷ ōṭavillai.} — Because of the strike the buses did not run today.
  \item \emph{Recall:} say the Tamil for a maid, then the Tamil for a trade union, then say \emph{vēlainiṟuttam} again
\end{itemize}

\begin{tcolorbox}[breakable,colback=teal!4,colframe=teal!35!black,title={Before you move on}]
What does \ta{வேலைநிறுத்தம்} mean? (\textquotedblleft{}A strike\textquotedblright{}.) Say it once more.
\end{tcolorbox}
\section[\texorpdfstring{\ta{பணிமனை} (paṇimaṉai)}{பணிமனை (paṇimaṉai)}]{\ta{பணிமனை} (paṇimaṉai) — a workshop; a depot}

\label{lesson:TA-C514-panimanai}



\begin{quote}
Before the new one: say the Tamil for a trade union, then the Tamil for a strike.
\end{quote}

\subsection*{You'll want to know: \ta{பணிமனை}}
\textbf{\ta{பணிமனை}} — \emph{paṇimaṉai} — \textquotedblleft{}a workshop; a depot\textquotedblright{}.

\begin{quote}
\ta{பழுதான} \ta{பேருந்தைப்} \ta{பணிமனைக்குக்} \ta{கொண்டு} \ta{போனார்கள்}.
\end{quote}

\begin{quote}
\emph{paḻutāṉa pēruntaip paṇimaṉaikkuk koṇṭu pōṉārkaḷ.}
\end{quote}

\begin{quote}
They took the broken-down bus to the depot.
\end{quote}

\ta{பணி} \textquotedblleft{}work\textquotedblright{} + \ta{மனை} \textquotedblleft{}a building, a house site\textquotedblright{}.

\begin{cousinweb}
\textbf{\ta{பணிமனை}} (\emph{paṇimaṉai}) is a native compound: \ta{பணி} (\emph{paṇi}), \textquotedblleft{}work\textquotedblright{} (DEDR 3884), and \ta{மனை} (\emph{maṉai}), \textquotedblleft{}a house, a building site\textquotedblright{}.
\end{cousinweb}

\begin{grammarlens}[title={What you've built}]
One more word for describing what you lived and felt.
\end{grammarlens}

\subsection*{Your turn}
\begin{itemize}
\raggedright
  \item \emph{Say it:} \emph{paṇimaṉai}
  \item \emph{Say it:} \emph{paṇimaṉai}, once more
  \item \emph{Say it:} \emph{paḻutāṉa pēruntaip paṇimaṉaikkuk koṇṭu pōṉārkaḷ.} — They took the broken-down bus to the depot.
  \item \emph{Recall:} say the Tamil for a trade union, then the Tamil for a strike, then say \emph{paṇimaṉai} again
\end{itemize}

\begin{tcolorbox}[breakable,colback=teal!4,colframe=teal!35!black,title={Before you move on}]
What does \ta{பணிமனை} mean? (\textquotedblleft{}A workshop; a depot\textquotedblright{}.) Say it once more.
\end{tcolorbox}
\section[\texorpdfstring{\ta{கிடங்கு} (kiṭaṅku)}{கிடங்கு (kiṭaṅku)}]{\ta{கிடங்கு} (kiṭaṅku) — a warehouse; a storeroom}

\label{lesson:TA-C514-kitanku}



\begin{quote}
Before the new one: say the Tamil for a strike, then the Tamil for a workshop.
\end{quote}

\subsection*{You'll want to know: \ta{கிடங்கு}}
\textbf{\ta{கிடங்கு}} — \emph{kiṭaṅku} — \textquotedblleft{}a warehouse; a storeroom\textquotedblright{}.

\begin{quote}
\ta{அரிசி} \ta{மூட்டைகள்} \ta{கிடங்கில்} \ta{இருக்கின்றன}.
\end{quote}

\begin{quote}
\emph{arici mūṭṭaikaḷ kiṭaṅkil irukkiṉṟaṉa.}
\end{quote}

\begin{quote}
The sacks of rice are in the warehouse.
\end{quote}

Where something is kept takes -\ta{இல்}: \ta{கிடங்கில்} \ta{வை}, to keep in the warehouse.

\begin{cousinweb}
\textbf{\ta{கிடங்கு}} (\emph{kiṭaṅku}) is native Dravidian: DEDR 1525 gives \ta{கிடங்கு} (\emph{kiṭaṅku}) for a storehouse.
\end{cousinweb}

\begin{grammarlens}[title={What you've built}]
One more word for describing what you lived and felt.
\end{grammarlens}

\subsection*{Your turn}
\begin{itemize}
\raggedright
  \item \emph{Say it:} \emph{kiṭaṅku}
  \item \emph{Say it:} \emph{kiṭaṅku}, once more
  \item \emph{Say it:} \emph{arici mūṭṭaikaḷ kiṭaṅkil irukkiṉṟaṉa.} — The sacks of rice are in the warehouse.
  \item \emph{Recall:} say the Tamil for a strike, then the Tamil for a workshop, then say \emph{kiṭaṅku} again
\end{itemize}

\begin{tcolorbox}[breakable,colback=teal!4,colframe=teal!35!black,title={Before you move on}]
What does \ta{கிடங்கு} mean? (\textquotedblleft{}A warehouse; a storeroom\textquotedblright{}.) Say it once more.
\end{tcolorbox}
\section[\texorpdfstring{\ta{பெறுநர்} (peṟunar)}{பெறுநர் (peṟunar)}]{\ta{பெறுநர்} (peṟunar) — a recipient; an addressee}

\label{lesson:TA-C514-perunar}



\begin{quote}
Before the new one: say the Tamil for a workshop, then the Tamil for a warehouse.
\end{quote}

\subsection*{You'll want to know: \ta{பெறுநர்}}
\textbf{\ta{பெறுநர்}} — \emph{peṟunar} — \textquotedblleft{}a recipient; an addressee\textquotedblright{}.

\begin{quote}
\ta{பொட்டலத்தில்} \ta{பெறுநரின்} \ta{முகவரி} \ta{தெளிவாக} \ta{இல்லை}.
\end{quote}

\begin{quote}
\emph{poṭṭalattil peṟunariṉ mukavari teḷivāka illai.}
\end{quote}

\begin{quote}
The recipient's address on the parcel is not clear.
\end{quote}

From \ta{பெறு} (\emph{peṟu}), to receive, with -\ta{நர்} \textquotedblleft{}one who …\textquotedblright{}; compare \ta{ஓட்டுநர்} (\emph{ōṭṭunar}).

\begin{cousinweb}
\textbf{\ta{பெறுநர்}} (\emph{peṟunar}) is native Dravidian: \ta{பெறு} (\emph{peṟu}), \textquotedblleft{}to get, to receive\textquotedblright{} (DEDR 4422), with -\ta{நர்}, \textquotedblleft{}one who does it\textquotedblright{}.
\end{cousinweb}

\begin{grammarlens}[title={What you've built}]
One more word for describing what you lived and felt.
\end{grammarlens}

\subsection*{Your turn}
\begin{itemize}
\raggedright
  \item \emph{Say it:} \emph{peṟunar}
  \item \emph{Say it:} \emph{peṟunar}, once more
  \item \emph{Say it:} \emph{poṭṭalattil peṟunariṉ mukavari teḷivāka illai.} — The recipient's address on the parcel is not clear.
  \item \emph{Recall:} say the Tamil for a workshop, then the Tamil for a warehouse, then say \emph{peṟunar} again
\end{itemize}

\begin{tcolorbox}[breakable,colback=teal!4,colframe=teal!35!black,title={Before you move on}]
What does \ta{பெறுநர்} mean? (\textquotedblleft{}A recipient; an addressee\textquotedblright{}.) Say it once more.
\end{tcolorbox}
